\documentclass[mnsc,nonblindrev]{informs3_NoOR} 
\OneAndAHalfSpacedXI
\usepackage{thmtools}
\usepackage[dvipsnames]{xcolor}
\usepackage{verbatim}
\usepackage{multirow}
\usepackage{float} 
\usepackage{endnotes}   
\usepackage{enumitem}   	 
\usepackage{amsmath}
\usepackage{mathrsfs}
\usepackage{mathtools}
\usepackage{url} 
\usepackage{subcaption} 
\usepackage[T1]{fontenc}
\usepackage[english]{babel}
\usepackage{lmodern} 
\usepackage{bbm}
\usepackage{lscape}
\usepackage{scalerel}
\usepackage{xr, xr-hyper}
\usepackage{hyperref}
\usepackage[colorinlistoftodos,prependcaption,textsize=tiny]{todonotes}
\usepackage{regexpatch}
\makeatletter
\xpatchcmd{\@todo}{\setkeys{todonotes}{#1}}{\setkeys{todonotes}{inline,#1}}{}{}
\makeatother
\hypersetup{
    colorlinks,
    linkcolor={red!50!black},
    citecolor={blue!50!black},
    urlcolor={blue!80!black}
}
\usepackage{natbib}
 \bibpunct[, ]{(}{)}{,}{a}{}{,}%
 \def\bibfont{\small}%
 \def\bibsep{\smallskipamount}%
 \def\bibhang{24pt}%
 \def\newblock{\ }%
 \def\BIBand{and}%
\TheoremsNumberedThrough 
\ECRepeatTheorems
\EquationsNumberedThrough   
\MANUSCRIPTNO{NoNumber}
\usepackage{soul}
\sethlcolor{yellow}

\begin{document}

\TITLE{Content Generation: Human Creators, AI Learning, and Platform Design}
\maketitle


\section{Introduction}

\section{Related Literature}
\begin{itemize}
\item \cite{ai_genai_2026}
\item \cite{esmaeili_how_2025}
\item \cite{taitler_braesss_2024}
\item \cite{taitler_data_2025}
\item \cite{taitler_selective_2025}
\item \cite{yao_rethinking_2023}
\item \cite{yao_human_2024}
\item \cite{yao_user_2024}
\item \cite{yu_beyond_2025}
\end{itemize}


%%%%%%%%%%%%%%%%%%%%%%%%%%%%%%%%%%%%%%%%%%%%%%%%%%%%%%%%%%
\section{Model}\label{sec:model}
%%%%%%%%%%%%%%%%%%%%%%%%%%%%%%%%%%%%%%%%%%%%%%%%%%%%%%%%%%
 
We analyze a platform model in which a human creator and an AI content generator compete for consumer attention. The platform sets strategic parameters in Stage~1; the human creator chooses effort in Stage~2. We solve by backward induction.
 
\subsection{Content Generation and Quality}
 
The human creator exerts effort $q_h \geq 0$ at cost $C(q_h) = \tfrac{1}{2}q_h^2$. The AI generates content of quality
\begin{equation}\label{eq:qA}
q_A = \alpha\, q_h + Q,
\end{equation}
where $\alpha \in (0,1]$ is the AI's learning efficiency and $Q \geq 0$ is its baseline quality.
 
\subsection{Consumer Demand}
 
Consumers are uniformly distributed on $[0,1]$, with location~$x$ indexing proximity to the AI (at~0) relative to the human (at~1). Consumer utility is
\begin{equation}\label{eq:utility}
u_A = q_A - t\,x\,m_A, \qquad u_h = q_h - t(1-x)\,m_H,
\end{equation}
where $t > 0$ is the mismatch cost, $\beta_h \in [0,1]$ is the platform's promotion weight for human content, $\delta \in (0,1)$ scales algorithmic influence, and
\begin{equation}\label{eq:mismatch}
m_A(\beta_h) = 1 - \delta(1-\beta_h), \qquad m_H(\beta_h) = 1 - \beta_h\delta.
\end{equation}
Both lie in $[1-\delta,\,1]$. Increasing $\beta_h$ raises $m_A$ (penalizing AI) and lowers $m_H$ (favoring human). The demand functions are
\begin{equation}\label{eq:demand}
D_A = \frac{q_A}{t\,m_A}, \qquad D_h = \frac{q_h }{t\,m_H}.
\end{equation}
 
\subsection{Platform and Creator Payoffs}
 
The platform incurs a discrimination penalty $\tfrac{k}{2}(\beta_h - \tfrac{1}{2})^2$, $k \geq 0$, representing the regulatory or reputational cost of departing from neutral curation. Under view-based revenue, the platform's payoff is
\begin{equation}\label{eq:up-view}
u_p^V = D_A + (1-r)\,D_h - \tfrac{k}{2}(\beta_h - \tfrac{1}{2})^2.
\end{equation}
Under engagement-based revenue it is
\begin{equation}\label{eq:up-eng}
u_p^E = q_A D_A + (q_h - r)\,D_h - \tfrac{k}{2}(\beta_h - \tfrac{1}{2})^2.
\end{equation}
The human creator earns
\begin{equation}\label{eq:uc}
u_c = r\,D_h - \tfrac{1}{2}q_h^2.
\end{equation}

%%%%%%%%%%%%%%%%%%%%%%%%%%%%%%%%%%%%%%%%%%%%%%%%%%%%%%%%%%
\section{Equilibrium Analysis: View-Based Revenue}\label{sec:view}
%%%%%%%%%%%%%%%%%%%%%%%%%%%%%%%%%%%%%%%%%%%%%%%%%%%%%%%%%%

\subsection{Stage 2: Creator's Optimal Effort}

\begin{proposition}\label{prop:effort}
Given a compensation rate $r \geq 0$, the equilibrium human effort is $q_h^* = \frac{r}{t m_H}$. Consequently, the induced AI content quality is $q_A^* = \frac{\alpha r}{t m_H} + Q$.
\end{proposition}

\subsection{Stage 1: Optimal Compensation}

\begin{proposition}\label{prop:rstar-view}
Under a view-based revenue model, the platform's payoff is strictly concave in the compensation rate~$r$. The unconstrained optimal compensation rate is:
\begin{equation}\label{eq:rhat-view}
\hat{r}^V = \frac{1}{2}\left(1 + \frac{\alpha m_H}{m_A}\right)
\end{equation}
Because $\hat{r}^V > \frac{1}{2}$, the platform always shares a strictly positive fraction of revenue to subsidize quality. 
\end{proposition}

\subsection{Stage 1: Optimal Promotion}


Let $W(\beta_h) = \frac{\alpha}{tm_A} + \frac{1}{tm_H} $, $\;y_i = 1/m_i$ for $i \in \{A,H\}$, and
\[
L(\beta_h) = \tfrac{1}{2}\,W\,W' + \tfrac{Q}{t}\,y_A' 
\]
denote the marginal revenue from promotion. Let $\Gamma(\beta_h) = L'(\beta_h)$ be the second derivative of the unpenalised utility function, and $\bar\Gamma = \max_{\beta_h \in [0,1]} \Gamma(\beta_h)$.
 
\begin{proposition}[Optimal Promotion: View-Based]\label{prop:beta-view}
The penalised utility function is
\[
U^V(\beta_h) = \tfrac{1}{4}W^2 + \tfrac{Q}{t}\,y_A  - \tfrac{k}{2}(\beta_h - \tfrac{1}{2})^2.
\]
The platform's optimal promotion weight is characterised as follows.
\begin{enumerate}
\item[\textup{(i)}] \textbf{$k = 0$:} $U^V$ is strictly convex and $\beta_h^* \in \{0,1\}$.
  
\item[\textup{(ii)}] \textbf{$k > \bar\Gamma$:} $U^V$ is strictly concave, so the maximum is unique. Define
\[
k^* = \max\!\big\{\bar\Gamma,\;\; {-2L(0)},\;\; 2L(1)\big\}.
\]
\begin{enumerate}
\item[\textup{(a)}] If $k > k^*$: the unique maximum is interior, $\beta_h^* \in (0,1)$, satisfying $L(\beta_h) = k(\beta_h - \tfrac{1}{2})$.
\item[\textup{(b)}] If $\bar\Gamma < k \leq k^*$ and $L(0) \leq -k/2$: the unique maximum is at $\beta_h^* = 0$.
\item[\textup{(c)}] If $\bar\Gamma < k \leq k^*$ and $L(1) \geq k/2$: the unique maximum is at $\beta_h^* = 1$.
\end{enumerate}
 
\item[\textup{(iii)}] \textbf{$0 < k \leq \bar\Gamma$:} $U^V$ may change concavity. The global maximum is found by comparing all stationary points of $L(\beta_h) = k(\beta_h - \frac{1}{2})$ with the boundary values $U^V(0)$ and $U^V(1)$.
\end{enumerate}
\end{proposition}

\begin{remark}\label{rem:k-view}
As $k \to 0$, $\beta_h^*$ converges to the corner solution; as $k \to \infty$, $\beta_h^* \to \tfrac{1}{2}$.
\end{remark}

%%%%%%%%%%%%%%%%%%%%%%%%%%%%%%%%%%%%%%%%%%%%%%%%%%%%%%%%%%
\section{Equilibrium Analysis: Engagement-Based Revenue}\label{sec:eng}
%%%%%%%%%%%%%%%%%%%%%%%%%%%%%%%%%%%%%%%%%%%%%%%%%%%%%%%%%%

\subsection{Stage 2: Creator's Optimal Effort}

The creator's problem is unchanged; Proposition~\ref{prop:effort} applies.

\subsection{Stage 1: Optimal Compensation}

\begin{proposition}\label{prop:rstar-eng}
Let $\Sigma(\beta_h) = \alpha^2 m_H + m_A(1 - tm_H ) $. 
\begin{enumerate}
\item[\textup{(a)}] \textbf{Concave Regime} ($\Sigma < 0$): The optimum is $r^* = \min\{\hat{r}^E, 1\}$, where:
\begin{equation}\label{eq:rhat-eng}
\hat{r}^E = \frac{\alpha t m_H^2Q }{-\Sigma}
\end{equation}
A sufficient condition for strict concavity ($\Sigma < 0$) globally is $t > \frac{1}{1-\delta} + \alpha^2$. 
\item[\textup{(b)}] \textbf{Convex/Linear Regime} ($\Sigma \geq  0$):The payoff is non-concave in $r$ and the platform's optimal subsidy is unbounded. The engagement formulation requires $\Sigma < 0$ for well-defined solutions.
\end{enumerate}
\end{proposition}

\subsection{Stage 1: Optimal Promotion}
Define the marginal promotion effects:
\begin{align}
\Omega_A(\beta_h) &= \frac{(2q_A^*)\,\alpha r^*\delta}{t^2 m_A m_H^2} - \frac{q_A^*(q_A^*)\,\delta}{tm_A^2}, \label{eq:OmA}\\[4pt]
\Omega_H(\beta_h) &= \frac{\delta}{tm_H^2}\!\left[\frac{r^*(2q_h^*-r^*)}{tm_H}+(q_h^*-r^*)(q_h^*)\right], \label{eq:OmH}
\end{align}
and let $F(\beta_h) = \Omega_A(\beta_h)+\Omega_H(\beta_h)$.
 
\begin{proposition}[Optimal Promotion: Engagement-Based]\label{prop:beta-eng}
Suppose $\Sigma = \alpha^2 m_H + m_A(1-tm_H) < 0$ at the equilibrium.
\begin{enumerate}
\item[\textup{(i)}] \textbf{$k=0$, $r^*>0$, $F(0)>0$:} For $Q$ sufficiently large, $\beta_h^* \in (0,1)$. For $Q$ small, $\beta_h^*=1$.
\item[\textup{(ii)}] \textbf{$k=0$, $r^*>0$, $F(0) \leq 0$:} $\beta_h^*=0$.
\begin{enumerate}
    \item If $F(\beta_h) \leq 0, \text{ for all } \beta_h \in [0,1]: \beta_h^*=0$.
    \item If $F(\beta_h) > 0, \text{ for some } \beta_h \in (0,1):$ the optimum is either $\beta^*_h=0$ or at an interior stationary point, determined by computing $U^E(0)$ with $U^E$ at each zero of $F$.
\end{enumerate}
\item[\textup{(iii)}] \textbf{$k=0$, $r^*=0$:} $\beta_h^*=0$.
\item[\textup{(iv)}] \textbf{$k>0$, $r^*>0$:} The optimum satisfies $F(\beta_h)=k(\beta_h-\frac{1}{2})$. The sign of $F(\frac{1}{2})$ determines the direction: $\beta_h^*>\frac{1}{2}$ if $F(\frac{1}{2})>0$, $\beta_h^*<\frac{1}{2}$ if $F(\frac{1}{2})<0$, and $\beta_h^*=\frac{1}{2}$ if $F(\frac{1}{2})=0$. As $k \to \infty$, $\beta_h^* \to \frac{1}{2}$.
\item[\textup{(v)}] \textbf{$k>0$, $r^*=0$:} $\beta_h^*<\frac{1}{2}$.
\end{enumerate}
\end{proposition}

% %%%%%%%%%%%%%%%%%%%%%%%%%%%%%%%%%%%%%%%%%%%%%%%%%%%%%%%%%%
% \section{Welfare Metrics}\label{sec:welfare}
% %%%%%%%%%%%%%%%%%%%%%%%%%%%%%%%%%%%%%%%%%%%%%%%%%%%%%%%%%%

% To compare outcomes across models and penalty levels, we use the following model-independent metrics:

% \begin{align}
% \text{Consumer surplus:}\quad & CS = \frac{(q_A^*)^2}{2tm_A} + \frac{(q_h^* )^2}{2tm_H}, \label{eq:cs}\\[4pt]
% \text{Creator welfare:}\quad & u_c^* = r^* D_h^* - \tfrac{1}{2}(q_h^*)^2, \label{eq:uc-metric}\\[4pt]
% \text{Normalized platform welfare:}\quad & u_p^* \label{eq:norm-up}
% \end{align}

% \[
% CS_A = \int_0^{D_A} \bigl(q_A  - t x m_A\bigr)\,dx
%       = q_A D_A - \frac{t m_A}{2}D_A^2
% \]
% Substituting \(D_A = \dfrac{q_A}{t m_A}\):
% \[
% CS_A
% = \frac{q_A^2}{t m_A}
%   - \frac{q_A^2}{2 t m_A}
% = \frac{q_A 2}{2 t m_A}
% \]
% Same logic for human content.
 

\bibliographystyle{plainnat}
\bibliography{references}

%%%%%%%%%%%%%%%%%%%%%%%%%%%%%%%%%%%%%%%%%%%%%%%%%%%%%%%%%%
\appendix
\section{Proofs}\label{app:proofs}
%%%%%%%%%%%%%%%%%%%%%%%%%%%%%%%%%%%%%%%%%%%%%%%%%%%%%%%%%%

\subsection*{Proof of Proposition~\ref{prop:effort}}
The human creator chooses $q_h$ to maximize:
\[
u_c = \frac{rq_h}{t m_H} - \frac{1}{2}q_h^2
\]
Differentiating with respect to $q_h$ yields the first-order condition:
\[
\frac{\partial u_c}{\partial q_h} = \frac{r}{t m_H} - q_h = 0 \implies q_h^* = \frac{r}{t m_H}
\]
The second derivative is $\partial^2 u_c/\partial q_h^2 = -1 < 0$, confirming a global maximum. Substituting this into the AI quality function $q_A = \alpha q_h + Q$ yields:
\[
q_A^* = \frac{\alpha r}{t m_H} + Q. \qquad \square
\]

\subsection*{Proof of Proposition~\ref{prop:rstar-view}}
Substitute the Stage-2 equilibrium qualities into the platform's view-based payoff (omitting the penalty term, which is independent of $r$):
\[
u_p^V = \frac{\frac{\alpha r}{t m_H} + Q }{t m_A} + (1-r)\frac{\frac{r}{t m_H} }{t m_H}
\]
Expanding and grouping by powers of $r$, we obtain a quadratic function $-Br^2 + Ar + C$, where:
\begin{equation}\label{eq:ABC}
B = \frac{1}{t^2 m_H^2}, \quad A = \frac{\alpha}{t^2 m_A m_H} + \frac{1}{t^2 m_H^2} , \quad C = \frac{Q }{t m_A}
\end{equation}
The second derivative with respect to $r$ is $-2B = -\frac{2}{t^2 m_H^2} < 0$, confirming strict concavity. The unconstrained maximum is $\hat{r}^V = \frac{A}{2B}$. By factoring $\frac{1}{t m_H}$ out of $A$, we get:
\[
A = \frac{1}{t m_H}\left(\frac{\alpha}{t m_A} + \frac{1}{t m_H} \right)
\]
\[
\hat{r}^V = \frac{A}{2B} = \frac{\frac{1}{t m_H}\left(\frac{\alpha}{t m_A} + \frac{1}{t m_H} \right)}{\frac{2}{t^2 m_H^2}} = \frac{1}{2}\left(\frac{\alpha m_H}{m_A} + 1 \right)
\]
Because all parameters are strictly positive, $\hat{r}^V > \frac{1}{2}$. \hfill$\square$

\subsection*{Proof of Proposition~\ref{prop:beta-view}}
\noindent\textbf{utility function derivation.}
Substituting $r^* = \hat{r}^V = A/(2B)$ into $u_p^V = -Br^2 + Ar + C - \frac{k}{2}(\beta_h-\frac{1}{2})^2$ gives $U^V = A^2/(4B) + C - \frac{k}{2}(\beta_h-\frac{1}{2})^2$. Using $A = \frac{1}{tm_H}\!\left(\frac{\alpha}{tm_A}+\frac{1}{tm_H}\right)$ and $B = \frac{1}{t^2 m_H^2}$:
\[
\frac{A^2}{4B} = \frac{1}{4}\!\left(\frac{\alpha}{tm_A}+\frac{1}{tm_H}\right)^{\!2} = \frac{1}{4}\,W^2.
\]
With $C = \frac{Q}{tm_A} = \frac{Q}{t}\,y_A $, the utility function is as stated.
 
\medskip\noindent\textbf{Preliminary facts.}
Since $m_A$ and $m_H$ are positive affine in $\beta_h$:
\[
y_A' = \frac{\delta}{m_A^2}>0,\quad y_A'' = \frac{2\delta^2}{m_A^3}>0;\qquad y_H' = \frac{\delta}{m_H^2}>0,\quad y_H'' = \frac{2\delta^2}{m_H^3}>0.
\]
$W = \frac{\alpha}{t}y_A + \frac{1}{t}y_H $ is strictly convex ($W''>0$), strictly positive, and strictly increasing ($W' = \frac{\alpha}{t}y_A'+\frac{1}{t}y_H'>0$). The equilibrium human quality is $q_h^* = \hat{r}^V/(tm_H) = \frac{1}{2}W$, and the participation constraint $q_h^*>0$ gives $W>0$.
 
The second derivative of $U^V$ is
\[
V''^V = \Gamma(\beta_h) - k, \qquad \Gamma = \tfrac{1}{2}(W')^2 + \tfrac{1}{2}WW'' + \tfrac{Q}{t}\,y_A'' .
\]
 
%---------- Part (i) ----------
\medskip\noindent\textbf{Part~(i): $k=0$.}
We show $\Gamma(\beta_h)>0$ for all $\beta_h \in [0,1]$.
 
\emph{Term~1:} $\frac{1}{2}(W')^2 \geq 0$.

\emph{Terms~2:} $\tfrac{1}{2}WW''$ since $W > 0$ and $W'' = \tfrac{\alpha}{t}\,y_A''+\tfrac{1}{t}\,y_H''>0$.
 
\emph{Term~3:} $\frac{Q}{t}\,y_A'' \geq 0$ since $Q \geq 0$ and $y_A''>0$.
 

All terms are non-negative and the last is strictly positive, so $\Gamma>0$. With $k=0$, $V''^V = \Gamma > 0$ everywhere: $U^V$ is strictly convex on $[0,1]$ and attains its maximum at a corner. \hfill$\square$
 
 
%---------- Part (ii) ----------
\medskip\noindent\textbf{Part~(ii): $k > \bar\Gamma$.}
Since $\Gamma$ is continuous on $[0,1]$, $\bar\Gamma<\infty$. For $k>\bar\Gamma$: $V''^V = \Gamma-k < 0$ everywhere, so $U^V$ is strictly concave and has a unique maximum.
 
The first derivative is $dU^V/d\beta_h = L(\beta_h) - k(\beta_h-\frac{1}{2})$. At the boundaries:
\[
\left.\frac{dU^V}{d\beta_h}\right|_0 = L(0)+\frac{k}{2}, \qquad \left.\frac{dU^V}{d\beta_h}\right|_1 = L(1)-\frac{k}{2}.
\]
 
\medskip\noindent\emph{Sub-case~(iii-a): $k > k^* \equiv \max\{\bar\Gamma,\,-2L(0),\,2L(1)\}$.}
Then $L(0)+k/2 > 0$ (derivative positive at $0$) and $L(1)-k/2<0$ (derivative negative at $1$). Since $U^V$ is strictly concave, increasing at the left boundary and decreasing at the right, the unique maximum is interior: $\beta_h^* \in (0,1)$, satisfying $L(\beta_h^*)=k(\beta_h^*-\frac{1}{2})$.
 
\medskip\noindent\emph{Sub-case~(iii-b): $\bar\Gamma < k \leq k^*$ with $L(0) \leq -k/2$.}
The derivative $dU^V/d\beta_h\big|_0 = L(0)+k/2 \leq 0$. Since $U^V$ is strictly concave and non-increasing at $\beta_h=0$, it is decreasing on all of $[0,1]$. The unique maximum is $\beta_h^*=0$.
 
This arises when promoting humans is sufficiently unprofitable at the left boundary (e.g., $Q$ large, $\alpha$ small) that even the penalty's push toward $\frac{1}{2}$ cannot overcome it.
 
\medskip\noindent\emph{Sub-case~(iii-c): $\bar\Gamma < k \leq k^*$ with $L(1) \geq k/2$.}
The derivative $dU^V/d\beta_h\big|_1 = L(1)-k/2 \geq 0$. Since $U^V$ is strictly concave and non-decreasing at $\beta_h=1$, it is increasing on all of $[0,1]$. The unique maximum is $\beta_h^*=1$.
 
This arises when human promotion is so valuable (e.g., $\alpha$ large, $Q$ small) that the penalty cannot restrain it. \hfill$\square$
 
%---------- Part (iii) ----------
\medskip\noindent\textbf{Part~(iii): $0 < k \leq \bar\Gamma$.}
The sign of $V''^V = \Gamma-k$ may vary: $U^V$ is locally convex where $\Gamma>k$ and locally concave where $\Gamma<k$. The utility function is continuous on $[0,1]$, so the Extreme Value Theorem guarantees a global maximum. It is found by:
\begin{enumerate}
\item solving $L(\beta_h) = k(\beta_h-\frac{1}{2})$ for all stationary points $\beta_h^{(1)},\ldots,\beta_h^{(n)} \in (0,1)$;
\item evaluating $U^V$ at each stationary point and at $\beta_h \in \{0,1\}$;
\item selecting the value yielding the largest $U^V$.
\end{enumerate}
This comparison is finite because the FOC is a smooth equation on a compact interval with generically finitely many solutions. \hfill$\square$

 
\subsection*{Proof of Proposition~\ref{prop:rstar-eng}}
Substitute $q_h^* = \frac{r}{t m_H}$ and $q_A^* = \frac{\alpha r}{t m_H} + Q$ into the engagement payoff:
\[
u_p^E = \frac{q_A^{*2} }{t m_A} + \frac{q_h^*(q_h^* - r) }{t m_H}
\]
\textbf{Step 1: AI Component Derivative.} Using $\frac{\partial q_A^*}{\partial r} = \frac{\alpha}{t m_H}$:
\[
\frac{\partial(q_A^* D_A)}{\partial r} = \frac{2q_A^*}{t m_A}\left(\frac{\alpha}{t m_H}\right) = \frac{2\alpha^2 r}{t^3 m_A m_H^2} + \frac{2Q\alpha}{t^2 m_A m_H}
\]
\textbf{Step 2: Human Component Derivative.} Let $H = \frac{(q_h^* - r)(q_h^*)}{t m_H}$. Substituting $q_h^*$:
\[
H = \left(\frac{r}{t m_H} - r\right)\frac{\frac{r}{t m_H} }{t m_H} = r\left(\frac{1 - t m_H}{t m_H}\right)\left(\frac{r }{t^2 m_H^2}\right) = \frac{1 - t m_H}{t^3 m_H^3}(r^2 )
\]
Differentiating with respect to $r$:
\[
\frac{\partial H}{\partial r}  = \frac{2r(1 - t m_H)}{t^3 m_H^3}
\]
\textbf{Step 3: First-Order Condition.} Summing the derivatives and multiplying the entire equation by $\frac{t^3 m_H^3 m_A}{2}$ to clear denominators yields:
\[
\alpha^2 m_H r + \frac{1}{2}\alpha t m_H^2 (2Q) + m_A(1 - t m_H)r= 0
\]
Factoring out $r$ and isolating the $\Sigma$ term:
\[
r\left[\alpha^2 m_H + m_A(1 - t m_H)\right] = -\alpha t m_H^2 Q
\]
Defining $\Sigma = \alpha^2 m_H + m_A(1 - t m_H)$ isolates $\hat{r}^E$ exactly as stated in the proposition. The second-order condition trivially follows the sign of $\Sigma$. \hfill$\square$

\subsection*{Proof of Proposition~\ref{prop:beta-eng}}
\noindent\textbf{Envelope derivative.}
Since $\partial u_p^E/\partial r = 0$ at $r^*$, the Envelope Theorem gives
\[
\frac{dU^E}{d\beta_h} = \frac{\partial}{\partial\beta_h}(q_A^* D_A) + \frac{\partial}{\partial\beta_h}((q_h^*-r^*)D_h) - k\!\left(\beta_h-\frac{1}{2}\right) = F(\beta_h) - k\!\left(\beta_h-\frac{1}{2}\right).
\]
The partial derivatives of the equilibrium qualities (holding $r$ fixed) are $\partial q_h^*/\partial\beta_h = r^*\delta/(tm_H^2)>0$ and $\partial q_A^*/\partial\beta_h = \alpha r^*\delta/(tm_H^2)>0$, using $m_H'=-\delta$ and $m_A'=\delta$.
 
\medskip\noindent\textbf{Derivation of $\Omega_A$.}
We have $q_A^* D_A = q_A^{*2}/(tm_A)$. Differentiating:
\[
\Omega_A = \frac{2q_A^*}{tm_A}\cdot\frac{\partial q_A^*}{\partial\beta_h} - \frac{q_A^{*2}}{tm_A^2}\cdot\delta = \frac{(2q_A^*)\alpha r^*\delta}{t^2 m_A m_H^2} - \frac{q_A^*(q_A^*)\delta}{tm_A^2}.
\]
The first term (positive effect) is positive when $q_A^*>0$ and grows linearly in $Q$. The second term (negative effect) is positive when $q_A^*>0$ and grows quadratically in $Q$ (through $q_A^{*2}$).
 
\medskip\noindent\textbf{Derivation of $\Omega_H$.}
We have $(q_h^*-r^*)D_h = (q_h^*-r^*)(q_h^*)/(tm_H)$. Differentiating using $\partial q_h^*/\partial\beta_h = r^*\delta/(tm_H^2)$:
\begin{align*}
\Omega_H &= \frac{\partial q_h^*}{\partial\beta_h}\cdot\frac{(2q_h^*-r^*)}{tm_H} + (q_h^*-r^*)(q_h^*)\cdot\frac{\delta}{tm_H^2} \\[4pt]
&= \frac{\delta}{tm_H^2}\!\left[\frac{r^*(2q_h^*-r^*)}{tm_H}+(q_h^*-r^*)(q_h^*)\right].
\end{align*}
Under $\Sigma<0$, $tm_H>1$, so $q_h^*=r^*/(tm_H)<r^*$ and $(q_h^*-r^*)<0$. The sign of $\Omega_H$ is therefore ambiguous and parameter-dependent.
 
%---------- Part (i) ----------
\medskip\noindent\textbf{Part~(i): $k=0$, $r^*>0$, $F(0)>0$.}
The derivative is $dU^E/d\beta_h = F(\beta_h)$.
 
\emph{At $\beta_h=0$:} $F(0)>0$ by assumption.
 
\emph{At $\beta_h=1$:} $m_A=1$, $m_H=1-\delta$, and $q_A^*(1) = \alpha r^*/(t(1-\delta))+Q$. The negative effect is $q_A^*(1)(q_A^*(1))\delta/t$. As $Q$ varies:
\begin{itemize}
\item The drag grows as $Q^2$: expanding $q_A^{*2}$, the leading term is $Q^2\delta/t$.
\item The positive effect grows as $Q$: the leading term is $2Q\alpha r^*\delta/(t^2(1-\delta)^2)$.
\item $\Omega_H$ is bounded independently of $Q$ (since $q_h^*$ does not depend on $Q$).
\end{itemize}
For $Q$ sufficiently large, the quadratic drag dominates and $F(1)<0$.
 
Since $F$ is continuous on $[0,1]$ with $F(0)>0$ and $F(1)<0$, the Intermediate Value Theorem yields at least one $\beta_h^* \in (0,1)$ where $F(\beta_h^*)=0$.
 
For $Q$ small, the drag is weak. If $F(\beta_h) \geq 0$ for all $\beta_h$, then $U^E$ is non-decreasing and $\beta_h^*=1$. \hfill$\square$
 
%---------- Part (ii) ----------
\medskip\noindent\textbf{Part~(ii): $k=0$, $r^*>0$, $F(0) \leq 0$.}
The derivative $dU^E/d\beta_h\big|_0 = F(0) \leq 0$: promoting human content is not marginally profitable even at the most favorable point.
 
If $F(\beta_h) \leq 0$ for all $\beta_h$ (which is the generic case when $F(0)<0$ strictly, since the negative effect only strengthens as $\beta_h$ grows), then $U^E$ is non-increasing and $\beta_h^*=0$: the platform fully promotes AI.
 
If $F$ changes sign in the interior (positive somewhere despite $F(0) \leq 0$), then $U^E$ may have an interior local maximum. Since $U^E$ is non-increasing near $\beta_h=0$, the global maximum is either at $\beta_h=0$ or at an interior stationary point, determined by direct comparison of $U^E$ values. In the generic case $F(0)<0$, $\beta_h^*=0$. \hfill$\square$
 
%---------- Part (iii) ----------
\medskip\noindent\textbf{Part~(iii): $k=0$, $r^*=0$.}
With $r^*=0$: $q_h^*=0$, $D_h=0$, and
\[
U^E(\beta_h) = \frac{Q^2}{t\,m_A(\beta_h)}.
\]
 
Since $m_A$ is increasing in $\beta_h$, $U^E$ is strictly decreasing. The maximum is at $\beta_h^*=0$, where $m_A=1-\delta$ is minimized and AI demand is maximized.
 
 
%---------- Part (iv) ----------
\medskip\noindent\textbf{Part~(iv): $k>0$, $r^*>0$.}
The derivative is $dU^E/d\beta_h = F(\beta_h)-k(\beta_h-\frac{1}{2})$. Define $G(\beta_h) = F(\beta_h)-k(\beta_h-\frac{1}{2})$.
 
\emph{At $\beta_h=\frac{1}{2}$:} The penalty term vanishes, so $G(\frac{1}{2})=F(\frac{1}{2})$.
\begin{itemize}
\item If $F(\frac{1}{2})>0$: $U^E$ is increasing at $\frac{1}{2}$, so $\beta_h^*>\frac{1}{2}$ (platform tilts toward human content).
\item If $F(\frac{1}{2})<0$: $U^E$ is decreasing at $\frac{1}{2}$, so $\beta_h^*<\frac{1}{2}$ (platform tilts toward AI).
\item If $F(\frac{1}{2})=0$: $\beta_h^*=\frac{1}{2}$ (neutral promotion).
\end{itemize}
 
\emph{At $\beta_h=0$:} $G(0)=F(0)+k/2$. The penalty contributes $+k/2>0$, boosting the derivative.
 
\emph{At $\beta_h=1$:} $G(1)=F(1)-k/2$. The penalty contributes $-k/2<0$, suppressing the derivative.
 
\emph{Convergence as $k \to \infty$:} Since $F$ is bounded on $[0,1]$ (let $M=\sup_{\beta_h}|F(\beta_h)|$), any zero of $G$ satisfies $|F(\beta_h^*)|=k|\beta_h^*-\frac{1}{2}|$, hence $|\beta_h^*-\frac{1}{2}| \leq M/k \to 0$ as $k \to \infty$. Therefore $\beta_h^* \to \frac{1}{2}$.
 
At the optimum, $F(\beta_h^*)=k(\beta_h^*-\frac{1}{2})$: marginal engagement gains equal marginal discrimination cost. \hfill$\square$
 
%---------- Part (v) ----------
\medskip\noindent\textbf{Part~(v): $k>0$, $r^*=0$.}
The payoff is $U^E = Q^2/(tm_A) - \frac{k}{2}(\beta_h-\frac{1}{2})^2$.
 
Let $\phi(\beta_h)=Q^2/(tm_A)$. Then $\phi'=-Q^2\delta/(tm_A^2)<0$: AI demand is decreasing in $\beta_h$. The FOC is
\[
\phi'(\beta_h) = k\!\left(\beta_h-\frac{1}{2}\right).
\]
The LHS is strictly negative for all $\beta_h$. The RHS is zero at $\frac{1}{2}$ and positive for $\beta_h>\frac{1}{2}$. Therefore equality requires $\beta_h^*-\frac{1}{2}<0$, i.e., $\beta_h^*<\frac{1}{2}$: the platform tilts toward AI because lower $\beta_h$ reduces $m_A$ and maximizes AI demand, while the penalty prevents going all the way to $\beta_h=0$.
 
The second derivative is $V''^E = \phi''-k$ where $\phi''=2Q^2\delta^2/(tm_A^3)>0$. For $k>\max_{\beta_h}\phi''$, $U^E$ is strictly concave and $\beta_h^*$ is the unique interior solution. For $k \leq \max_{\beta_h}\phi''$, the optimum may be at $\beta_h=0$ and is found by comparing the boundary with any interior stationary point. \hfill$\square$

\end{document}